\lesson{6}{Oct 14 2021 Thu (06:00:23)}{Periodic Table}{Unit 2}

There are many ways to sort items. Often, the easiest way to organize things is
by their similarities. All items could be grouped with other items of the same

\bf{color} or \bf{the same type}, such as:

\begin{itemize}
    \item Plants
    \item Animals
\end{itemize}

To group elements, \bf{Mendeleev} compared their:

\begin{itemize}
    \item Physical properties 
    \item Chemical properties 
\end{itemize}

and searched for common traits.

\subsubsection*{History of the Periodic Tabel}

It's been more than $200$ years since the Periodic Table of elements was first
described by Dmitri Mendeleev, but the idea of an element as a basic type of
material must have occurred to early humans as they selected certain types of
rocks for tools and later melted metal into weapons and baked wheat into bread.

Greek scholar \bf{Democritus} thought that substances were made up of tiny invisible
particles that could not be broken into smaller pieces.

Well, atoms—"atomos" in Greek means indivisible. And the way he defined it was
by imagining, say, a gold wire that you rolled incredibly thin, and you chopped
it in half and chopped it in half again, and then again, and then again.
Eventually you would get to something so small that you couldn't chop it in half
any more, and it would still be gold. That is an indivisible unit, an atom of
gold.

Other ancient Greek scholars classified the fundamental elements as:

\begin{itemize}
    \item Fire
    \item Earth
    \item Air
    \item Water
\end{itemize}

The philosopher Aristotle believed that these elements could be
mixed to create new materials. This sparked a centuries-long interest in
alchemy, specifically transmuting ordinary metals like lead into gold.

Scottish alchemist Robert Boyle concluded in the late $1600$s, however, that
although metals like gold and lead are both made of atoms, the size and shape of
the atoms are different.

Nearly $100$ years later, French chemist Antoine Laurent de Lavoisier offered the
first modern definition of an element, a chemical substance that could not be
broken down into another substance. He named more than $30$ elements himself.

His contemporary, John Dalton, figured out a way to estimate the weight of each
kind of atom. Hydrogen always made up 15 percent of the weight of water, whereas
oxygen accounted for $85$ percent of water's weight. Using these numbers, Dalton
was able to calculate the atomic weight of these two elements.

The list of elements and their atomic weights grew throughout the $18th$ and early
$19th$ centuries. But it remained just a list arranged from lightest to heaviest.
However, some scientists started to recognize patterns within the list.

There are a couple of patterns detected in the lists of elements. The first are
triads, where you take three elements in a row. The atomic weight of the middle
element would be the average of the atomic weight of the lighter and the
heavier.

Similarly, other chemists noticed that there were natural families in which
certain elements grouped. The halogens, for example, which are salt-making
elements, usually very highly reactive, tended to have very similar properties,
form very similar crystals, and so on. And those were grouped together as a
natural family.

Dmitri Mendeleev was writing a new textbook for his students when he noticed a
repeating, or periodic, pattern among the $63$ known elements.

His periodic table made its debut in $1869$. His table looks a little different
from the ones that are hanging in chemistry classrooms. The atomic weight
increases from top to bottom, and the families run from left to right. So you
would have to rotate it clockwise by $90$ degrees and then reflect it in the
mirror to make it look like our current one.

But it still has the basic properties of gradual increase of atomic weight and
periodic repetition of properties across families as you increase the weight.
Mendeleev put question marks in spaces where he thought there should be an
element of a certain atomic weight. Even though no element of that weight had
been discovered yet, he predicted their atomic weights and other features such
as their melting and boiling points.

In this sense, Mendeleev did not invent the periodic table of elements. He
discovered it. His carefully stacked rows and columns turned a simple list into
a snapshot of how matter is organized on Earth and throughout the universe.

Mendeleev looked for similarities between elements to group them together. He
started by grouping elements that were metals and nonmetals based on their
physical properties.

\subsubsection*{Physical Properties of Metals:}

\begin{itemize}
    \item Good conductors of electricity and heat
    \item Solid at room temperature
    \item Malleable, flexible, and ductile
    \item Lustrous (shiny)
    \item Higher density
\end{itemize}

\subsubsection*{Physical Properties of Nonmetals:}

\begin{itemize}
    \item Poor conductors of electricity and heat
    \item Solid, liquid, or gas at room temperature
    \item Solids are brittle and break easily
    \item Dull (not shiny)
    \item Lower density
\end{itemize}

If an element did not fully fit either grouping, it went to a group called
metalloids, which had characteristics of both metals and nonmetals.

Reading the periodic table takes a little practice. The letter or letters that
symbolize the element are sometimes self-evident, like $H$ for hydrogen. Some are
not so obvious, such as $Sb$ for antimony.

The atomic number at the top of the box represents the number of protons
contained in an atom of that element. Underneath the element's name is the
atomic weight. Each row of the table is called a period.

All of the elements in a period have the same number of electron shells. Columns
in the table are called groups. For most groups the elements have the same
number of electrons in their outermost electron shell.

Another way to look at the table is to divide the elements into metals,
nonmetals, and metalloids. Most of the elements in the table are metals. When
people are asked to name any element, the chances are good that they will come
up with the name of a transition metal.

Transition metals are found in numerous products such as:

\begin{itemize}
    \item Coins
    \item Jewelry
    \item Light
    \item Bulbs
    \item Cars
\end{itemize}

Transition metals are called the bridge of the periodic table and include groups
$3$ through $12$.

\it{Group 1 elements}, beginning with \it{lithium} and running
vertically to \it{francium}, are called \it{alkali metals}.

Group 2 elements, beginning with beryllium and running vertically to radium, are
called alkaline earth metals. Elements in these first two groups of the periodic
table are some of the most reactive elements known.

In fact, many of the elements are so reactive that water, or in some cases air,
can cause them to explode or catch fire. Although many of their names are
familiar:

\begin{itemize}
    \item Sodium
    \item Potassium
    \item Calcium
\end{itemize}

for instance—it is rare to find these elements on their own. Instead, these
elements are usually combined with another, more stable element to create a
compound.

The elements in \it{groups $13$} through $16$ of the periodic table are a
mixed bag. Some of the elements are metals, though they barely act like metals
in chemical reactions.

Other elements within these groups are definitely nonmetals, whereas still other
elements—metalloids—have a combination of metal and nonmetal features.
Scientists find it difficult to agree on one name for all these groups, so many
researchers call them the BCNOs after the first element in each group—boron,
carbon, nitrogen, and oxygen. BCNOs, depending on whether they act more like a
metal or a nonmetal, may give up electrons or share electrons with other atoms.
The valence shells of the BCNOs are about half full.

\begin{definition}[Atomic Mass]
    The atomic mass, which is measured in atomic mass units (amu), is an average of the masses of all the isotopes of that element.

    Most of the mass of an atom comes from the protons and neutrons. The electrons are so small they don't add much mass to the total.
\end{definition}

\begin{definition}[Atomic Number]
    Elements of the periodic table are arranged in order by their atomic
    numbers. The atomic number, which is usually found above the element's
    chemical symbol, indicates the number of protons in a neutral atom. Since
    positive charges equal negative charges in a neutral atom, the number of
    protons equals the number of electrons.
\end{definition}

\begin{definition}[Atomic Symbol]
    Chemical symbols are shorthand versions of the full names of elements. For
    each element, the chemical symbol is usually derived from the first letter
    or two of the element name. However, in some instances, it does not follow
    that pattern.
\end{definition}

The number of:

\begin{itemize}
    \item Protons
    \item Neutrons
    \item Electrons
\end{itemize}

within an atom affects its properties. You can determine the number of each of
these subatomic particles when given the atomic number and mass number for a
given atom.

\subsubsection*{Calculating Average Atomic Mass}

\paragraph{Differeing Neutrons} Isotopes of an element have the same number of
protons but a different number of neutrons. This means isotopes of the same
element will have different atomic masses.

Carbon has three isotopes: carbon-12, carbon-13, and carbon-14. All three
isotopes have six protons, but their neutron numbers are six, seven, and eight.
This means that the three isotopes have different atomic masses but share the
same atomic number of six.

The number added to the element's name refers to the mass number. So, carbon-13
is an isotope of carbon with six protons and seven neutrons, and carbon-14 is an
isotope of carbon with six protons and eight neutrons.

\paragraph{Differeing Masses} Most elements found in nature occur as a mixture
of two or more isotopes. However, the percentage of each isotope in nature is
different. Some isotopes are more common than others.

To calculate the average atomic mass of an element, the percent abundance of
each isotope of that element is multiplied by its mass. Then all the atomic
masses are added together to create a weighted average that accounts for the
percentage of each isotope in nature.

(percent abundance × mass) + (continue for each isotope) = average atomic mass


(Percent Abundance * mass) + Continue for each Isotop = Average Atomic Mass

\paragraph{Isotope Reps} There are two main ways to represent an isotope. Simply
place a dash to the right of the name with the mass number, such as uranium-235.

Or, place the mass number as a superscript to the upper left of the chemical
symbol, and place its atomic number as a subscript to the lower left of the
chemical symbol.

\newpage
